% Generated draft: requires mathematical and source review.
\documentclass[11pt]{article}
\usepackage[margin=25mm]{geometry}
\usepackage{fontspec}
\setmainfont{Latin Modern Roman}
\usepackage{amsmath,amssymb,mathtools,unicode-math}
\setmathfont{Latin Modern Math}
\usepackage{enumitem,needspace,xurl,hyperref}
\setlength{\parindent}{0pt}
\setlength{\parskip}{5pt}
\setlength{\emergencystretch}{3em}
\newcommand{\N}{\mathbb{N}}
\newcommand{\Z}{\mathbb{Z}}
\newcommand{\Q}{\mathbb{Q}}
\newcommand{\R}{\mathbb{R}}
\newcommand{\C}{\mathbb{C}}
\newcommand{\F}{\mathbb{F}}
\newcommand{\A}{\mathbb{A}}
\newcommand{\PP}{\mathbb{P}}
\newcommand{\E}{\mathbb{E}}
\newcommand{\eps}{\varepsilon}
\newcommand{\dd}{\,\mathrm{d}}
\begin{document}
\section*{UnsolvedMath ID: 20001150}
\textbf{Source problem code:} catalogue entry
\subsection{Definitions and mathematical statement}
A hole in a finite simple graph is an induced cycle of length at least four; a graph is even-hole-free if it has no hole of even length. Write $\omega(H)$ for clique number. A $k$-division of a graph $H$ with an edge is a partition into $k$ possibly empty sets $V_1,\ldots,V_k$ such that $\omega(H[V_i])<\omega(H)$ for every $i$. A graph is $k$-divisible if every induced subgraph with at least one edge has a $k$-division. The edge exception is essential because singleton maximum cliques cannot be nonmonochromatic.

Hoang's source conjecture is
\[
 \forall G:\quad G\text{ even-hole-free}\ \Longrightarrow\ G\text{ is $3$-divisible}.
\]
Equivalently every induced $H\subseteq G$ with an edge has a three-color assignment in which no maximum-cardinality clique is monochromatic. This concerns maximum cliques, not all inclusionwise maximal cliques.

The aid attributes to Chudnovsky--Seymour (2023) the existence of a bisimplicial vertex in every nonempty even-hole-free graph, and to Chen--Lan--Zhong a May 2026 preprint explicitly confirming the conjecture. Bisimplicial means that the neighborhood is the union of two cliques, giving $d_H(v)\leq2(\omega(H)-1)$.

Its stronger candidate is: for any integer $q\geq1$, if $G$ is a graph such that every induced subgraph $J\subseteq G$ with an edge has a vertex $v$ with
\[
 d_J(v)\leq q(\omega(J)-1),
\]
then for every induced $H$ with an edge and every assignment of lists $L(v)$ of $q+1$ distinct colors to its vertices, there is $c(v)\in L(v)$ such that no maximum clique of $H$ is monochromatic. In particular it reports hereditary three-list-colorability of the maximum-clique hypergraph for even-hole-free graphs. The lists need not share one common palette. The graph $C_5$ is even-hole-free with clique number two and cannot properly be two-colored, showing that three cannot be replaced by two for this class. These are supplied literature and research-aid conclusions requiring independent verification.
\subsection{Short English statement}
Every induced subgraph with an edge should admit three parts, each missing all maximum cliques. The supplied aid reports Hoang's conjecture solved and a stronger version permitting arbitrary three-element color lists; it does not assert the analogous condition for every maximal clique.
\subsection{Sources}
\begin{itemize}
\item AIM workshop \emph{The Perfect Graph Conjecture}, version August 24, 2004, Chapter G.4, Conjecture 2 attributed to Hoang: \url{https://aimath.org/WWN/perfectgraph/perfectgraph.pdf}.
\item Local registry \texttt{data/unsolvedmath/problems.json}, permanent ID 20001150, exact source code \texttt{catalogue entry}; canonical location \texttt{aim-combinatorics-notes.json notes[274]}. Companion ID 20001149 (\texttt{AIM-COMBINATORICS-0274}, \texttt{notes[273]}) contains the preceding definition, including the induced-subgraph edge exception.
\item The attached aid names Chudnovsky--Seymour (2023) and Chen--Lan--Zhong (May 2026 preprint), without full references. Attempt 1, full solution, problem status at run solved, rigor proved, low novelty confidence; machine-generated, not peer reviewed, requiring independent verification before citation.
\end{itemize}
\end{document}
